\textbf{A model-selection pitfall and its fix.} Our first MoE run (v1) selected checkpoints with $S$ computed from the \emph{mean} PSNR of the validation set. Because the identity branch returns clean images almost unchanged, their PSNR is $\approx$100\,dB, and these few values dominated the mean: the validation mean PSNR fell from 39.3 to $\approx$36.8\,dB as soon as joint fine-tuning started, although SSIM \emph{rose} from 0.865 to 0.871. The criterion therefore kept the end of the warm-up (experts still frozen) as the ``best'' model, the v1 MoE was in effect hard routing with a re-trained gate (corrupted inputs 26.17\,dB / 0.811 SSIM, routing entropy 0.07, 97\,\% of inputs with $\max_k w_k>0.9$), and the joint fine-tuning had been discarded. We found this by inspecting the per-epoch history rather than only the final numbers. The fix caps PSNR at 40\,dB \emph{per image} before averaging; the whole MoE stage (Optuna search and final training) was then re-run (v2) from the same Task-2 checkpoints. All MoE results below are from v2; the v1 study is kept in the repository as evidence.

\textbf{Reconstruction.} With the corrected criterion the best checkpoint is the last joint epoch (Fig.~\ref{fig:curves}, bottom): validation SSIM rises from 0.871 after warm-up to 0.884, so joint fine-tuning now clearly helps. On the test set the soft MoE is the best system on corrupted inputs, with \MOEPSNRCORR\,dB / \MOESSIMCORR{} SSIM versus \TONEPSNRCORR\,dB / \TONESSIMCORR{} (universal DAE) and \PREDPSNRCORR\,dB / \PREDSSIMCORR{} (predicted hard routing). It is best in SSIM for every corruption and severity (Table~\ref{tab:ssim_sev}) and best in PSNR for all blur and salt-and-pepper levels. The largest gains appear exactly where the bottleneck hurt the other systems: low blur (29.7\,dB / 0.901 vs.\ 28.1 / 0.857 for hard routing) and low occlusion (SSIM 0.881 vs.\ 0.805). Clean images keep near-identity quality (58.5\,dB, SSIM 0.999). The only regression is occlusion PSNR at medium/high severity (20.56 vs.\ 20.77\,dB at high severity), discussed below.

\begin{figure}[t]
\centering
\includegraphics[width=0.85\columnwidth]{t3_routing_heatmap.png}
\caption{Routing heatmap: mean gate weight of each branch for every true corruption and severity on the test set.}
\label{fig:heat}
\end{figure}

\begin{figure*}[t]
\centering
\includegraphics[width=0.98\textwidth]{t3_weight_distribution.png}
\caption{Distribution of the four routing weights per true input condition (test set, box = inter-quartile range).}
\label{fig:wdist}
\end{figure*}

\textbf{Routing behaviour} (Figs.~\ref{fig:heat}, \ref{fig:wdist}). The gate remains semantically aligned with the corruption: the matching branch receives the largest weight for clean (0.96), salt-and-pepper (0.89--0.96), blur (0.65--0.87) and occlusion (0.56--0.73) inputs, and the argmax of $w$ agrees with the true label for 96.7\,\% of test inputs. No expert is inactive (overall mean weights 0.27/0.28/0.25/0.20), and no expert dominates unrelated inputs: the salt-and-pepper, blur and occlusion experts receive on average only 0.0002, 0.019 and 0.008 of the weight on inputs that are not of their type. The only branch used on unrelated inputs is the \emph{identity} branch (0.19 on average for corrupted inputs) --- and this is the key learned behaviour. The identity weight decreases monotonically with severity (blur: 0.33/0.16/0.13; occlusion: 0.44/0.32/0.27; salt-and-pepper: 0.06/0.02/0.01): for mild corruptions the gate blends a sizeable fraction of the (almost clean) input with the expert's smoothed reconstruction, which restores the high-frequency detail that the bottleneck removes. This is effectively a learned, input-adaptive skip connection that hard routing cannot express. It already appears when only the gate is trained with the higher temperature of v2 (mean validation identity weight 0.32 after warm-up, where 0.25 would correspond to clean inputs only) and is strengthened by joint fine-tuning (0.38, Fig.~\ref{fig:curves}), during which the experts adapt to being blended. The routing entropy increased from 0.07 (v1) to 0.45, and only 35\,\% of inputs now have one branch with $w>0.9$.

\begin{figure}[t]
\centering
\includegraphics[width=\columnwidth]{t3_dominant_examples.png}\\[1mm]
\includegraphics[width=\columnwidth]{t3_distributed_examples.png}
\caption{Top: inputs where one expert dominates ($w>0.9$). Bottom: inputs with the most distributed weights ($w=[w_0,w_1,w_2,w_3]$ in each title).}
\label{fig:moeex}
\end{figure}

\textbf{Dominant vs.\ distributed examples} (Fig.~\ref{fig:moeex}). Strong salt-and-pepper noise is handled almost exclusively by the salt-and-pepper expert ($w_1\approx0.95$; the blur expert receives a small, consistent 0.03--0.05). The most distributed weights occur for \emph{low}-severity salt-and-pepper on highly textured images (grass, flowers), where the gate splits $\approx$0.45/0.50 between identity and the salt-and-pepper expert: the texture makes the sparse impulses hard to distinguish from image detail. This is also a failure mode of soft routing --- blending 45\,\% of the input re-introduces faint impulses, visible as dots in the error maps. The same mechanism explains the occlusion regression: blending 27--44\,\% identity darkens the in-painted region again, which lowers PSNR inside large masks although SSIM, dominated by the unoccluded area, improves.

\textbf{Mixed corruptions.} To test the motivation of soft routing, we applied blur $(5,1.5)$ followed by salt-and-pepper $p=0.08$ --- a combination never seen in training --- to 300 test images. The classifier labels all of them as salt-and-pepper and the gate also concentrates on that expert ($w_1=0.95$, $w_2=0.04$). The MoE (25.8\,dB / 0.776) is only marginally better than hard routing (25.7\,dB / 0.771), while the universal DAE is clearly best (27.0\,dB / 0.821). Soft routing alone therefore does not solve compound degradations: the gate was never shown such inputs and the experts are applied in parallel, not in sequence, so no weighted sum of single-corruption experts can undo both corruptions. The universal model, which shares one representation across corruptions, generalises better here.

\textbf{Optuna.} Three of the eight v2 trials completed and five were pruned by the median rule; none had to be stopped by the routing-collapse rule, because the classifier initialisation and the cross-entropy term keep the gate aligned from the start. The search separates cleanly by temperature: all three completed trials used $\tau\ge1.6$ and all five pruned trials $\tau\le1.29$. The selected configuration (\MOEBESTPARAMS) uses a learning rate near the top of the range, a high temperature (softer weights) and small classification and balance weights, which is consistent with the observed benefit of softer, blended routing.
